\section{数据系统：质量、许可与“数据污染”三重约束}

\subsection{数据混配与合成数据成为常规能力}

对话中提到的数据工程实践非常务实：先做小样本混配实验，再根据目标评测回归最佳配比。这个流程反映的是“任务驱动数据”而非“数据越多越好”。同时 OCR 与高质量合成样本让可用 token 继续增长。

\begin{importantbox}{合成数据的正确打开方式}
合成数据真正有价值的前提是“可验证过滤”：
\begin{itemize}[leftmargin=1.5em]
    \item 数学题：程序化验算；
    \item 代码题：单元测试与执行验证；
    \item 复杂问答：多模型交叉一致性 + 人工抽检。
\end{itemize}
这意味着“生成”只是开始，关键在后续的筛选、标注与版本管理。
\end{importantbox}

\subsection{法律边界与数据主权}

访谈触及了真实法律风险：训练语料是否有授权、历史抓取是否合规、版权责任如何追溯。这个问题不会被技术进步自动抹平。越到后期，数据主权越会变成组织护城河：谁能合法获得高质量私有数据，谁就更有机会做出同质化之外的能力。

\begin{warningbox}{“抓得到就能训”是高风险思路}
在高压监管环境下，数据来源证明链将越来越像财务审计。没有可追溯来源的语料，短期可能加速训练，长期却可能引发高额赔偿、模型下架与品牌损失，风险远超一次训练收益。
\end{warningbox}

\subsection{LLM 数据污染与开源维护者负担}

随着网络内容中 AI 生成比例上升，训练语料开始出现“模型学模型”的回环风险。Sebastian 对开源维护社区的观察很实际：大量低质量 AI 辅助 PR 正在消耗维护者精力。这里的关键不是“AI 生成”本身，而是是否有人类验证层。

\begin{knowledgebox}{判断“可用数据”的一个朴素标准}
如果一段内容无法被执行、对照、复现或交叉核验，那么它很难成为高价值训练样本。代码语料相对好处理，因为可运行；概念阐释与评论文本更难评估质量，需要额外评价管线。
\end{knowledgebox}

\subsection{本章小结}

2026 的数据工程已经进入“数据质量管理学”阶段：混配策略、法律合规、污染治理同等重要，缺一不可。

